\chapter{结论与展望}

\section{基本结论}

习题课是不可或缺的高中数学课型之一，探究性学习是高效学习的良好方法，因此高中数学习题课中的探究性学习需要被发现、被重视。但在习题课下的探究性学习与一般的探究性学习有较大区别，因此笔者通过整理前人的研究进行文献综述，结合学者的相关概念定义，给出了数学习题课中的探究性学习的定义，该定义强调了常规习题的基础性与习题延伸的可探究性，用探究性学习的一般步骤进一步获得习题与其延伸的解决方法与路径。

本文根据理论研究给出基于探究性学习的习题课教学实例方案：函数的凹凸性的探究性习题教学，在两个班级进行实验后进行实施效果分析，得到学生与教师两个层面的反馈。从学生层面反馈上看，教学三维目标的实现较为良好，其中在掌握研究路径这一方面得到良好反馈；从教师层面反馈上看，多数旁听教师认为学生主动性和课堂积极性的提升是该实例的重要成效。但该实例也存在一定不足，在小组交流、时间掌握方面还有待改善。



\section{研究展望}

在该课题的研究过程中笔者也遇到了一些困难，例如探究性习题课的开展对班级优等生、中等生以及后进生的影响有何偏差？在实例调查的过程中，有学生向笔者反应自己平时学习成绩跟不上，在探究性课堂也难以跟进，笔者认为是否可以将学生进行分类教学，不同层次的学生进行不同层次的习题探究？笔者在研究教学实例时发现，并不是所有习题都适合进行探究性学习，因此什么样的习题才适合发展探究性学习呢？在今后的教学工作中，笔者将会继续探索高中数学中的探究性学习，不断地完善与改进数学探究学习方法。